% concatenated from 241101339.tex

\documentclass[11pt]{amsart}%
\usepackage{amscd,amsmath,latexsym,amsthm,amsfonts,amssymb,graphicx,color,geometry,hyperref}
\newcommand {\fkf} {\mathfrak{F}}
\usepackage[utf8]{inputenc}
\usepackage{amsmath}
\usepackage{amsfonts}
\usepackage{amssymb}
\usepackage{graphicx}
\setcounter{MaxMatrixCols}{30}
\usepackage[all]{xy}
\usepackage{indentfirst}
\usepackage{graphicx}
\usepackage{graphics}
\usepackage{pict2e}
\usepackage{epic}
\usepackage{array} 
\numberwithin{equation}{section}
\usepackage{dsfont}
\providecommand{\U}[1]{\protect\rule{.1in}{.1in}}
\vfuzz2pt
\hfuzz2pt
\theoremstyle{plain}
\newtheorem{thm}{Theorem}[section]
\newtheorem{lem}[thm]{Lemma}
\newtheorem{cor}[thm]{Corollary}
\newtheorem{dfn}[thm]{Definition}
\newtheorem{prop}[thm]{Proposition}
\theoremstyle{definition}
\newtheorem{defn}[thm]{Definition}
\newtheorem{Conj}[thm]{Conjecture}
\newtheorem{rem}[thm]{Remark}
\newtheorem{?}[thm]{Problem}
\newtheorem{exa}[thm]{Example}
\theoremstyle{definition}
\newtheorem*{nt*}{Notation}
\newcommand{\norm}[1]{\left\Vert#1\right\Vert}
\newcommand{\abs}[1]{\left\vert#1\right\vert}
\newcommand{\set}[1]{\left\{#1\right\}}
\newcommand{\inner}[2]{\left\langle {#1}, {#2} \right\rangle}


\def\dlim{\displaystyle\lim}

\newcommand{\im}{{\rm Im}\,}
\newcommand{\re}{{\rm Re}\,}
\newcommand{\dist}{\hbox{ \rm dist}}
\newcommand{\Span}{\hbox{Span}}

\newcommand{\calC}{\mathcal C}
\newcommand{\calE}{\mathcal E}
\newcommand{\calF}{\mathcal F}
\newcommand{\calO}{\mathcal O}
\newcommand{\calQ}{\mathcal Q}
\newcommand{\calJ}{\mathcal J}
\newcommand{\calK}{\mathcal K}

\newcommand {\C} {\mathbb C}
\newcommand {\R} {\mathbb R}
\newcommand {\Z} {\mathbb Z}
\newcommand {\X} {\mathbb X}
\newcommand {\V} {\mathbb V}



\def\dlim{\displaystyle\lim}
\def\dsum{\displaystyle\sum}

\begin{document}
\title[]{Cyclicity of composition operators on the Paley-Wiener spaces}
	
	
\author{Pham Viet Hai, Waleed Noor and Osmar Reis Severiano}%
\address[P. V. Hai]{Faculty of Mathematics and Informatics, Hanoi University of Science and Technology, Khoa Toan-Tin, Dai hoc Bach khoa Hanoi, 1 Dai Co Viet, Hanoi, Vietnam.}%
\email{hai.phamviet@hust.edu.vn}
\address[W. Noor]{IMECC, Universidade Estadual de Campinas, Campinas-SP, Brazil}
\email{waleed@unicamp.br}
\address [O. R. Severiano]{ IMECC, Universidade Estadual de Campinas, Campinas, Brazil}
\email{osmar.rrseveriano@gmail.com}
\thanks{P.V. Hai is funded by Vietnam National Foundation for Science and Technology Development (NAFOSTED) under grant number 101.02-2021.24, O.R. Severiano is a postdoctoral fellow at the Programa de Matemática
and is supported by UNICAMP (Programa de Pesquisador de Pós-Doutorado PPPD) }
	
	
	
	
\keywords{Cyclic operator; composition operator; Paley-Wiener space}
	
\begin{abstract} In this article we characterize the cyclicity of bounded composition operators $C_\phi f=f\circ \phi$ on the Paley-Wiener spaces of entire functions $B^2_\sigma$ for $\sigma>0$. We show that $C_\phi$ is cyclic precisely when $\phi(z)=z+b$ where either $b\in\mathbb{C}\setminus\mathbb{R}$ or $b\in\mathbb{R}$ with $0<|b|\leq \pi/\sigma$. We also describe when the reproducing kernels of $B^2_\sigma$ are cyclic vectors for $C_\phi$ and see that this is related to a question of completeness of exponential sequences in $L^2[-\sigma,\sigma]$. The interplay between cyclicity and complex symmetry plays a key role in this work.
\end{abstract}
	
\maketitle
	
	
	
%===================================================================================================
\section{Introduction}
%===================================================================================================
% Let $\mathbb{C}$ denote the complex plane and $\mathbb{R}$ the real line. An entire function $f$ is said to be of \textit{exponential type} if the inequality $
%|f(z)|\leq Ae^{B|z|}$ holds for all $z\in \mathbb{C}$
%and for some positive constants $A$ e $B.$ In this case, the \textit{exponential type} of $f$ is defined to be the number 
%\begin{align*}
%%\sigma=\limsup\limits_{r\rightarrow\infty}\frac{\log M_f(r)}{r},
%\end{align*}
%where $M_f(r)=\max\{|f(z)|:|z|=r\}. $ The identically zero function has exponential type zero, by convention.  We refer the reader to \cite{Branges,Levin,young} for background on entire functions. For $\sigma>0,$ the \emph{Paley-Wiener space} $B_\sigma^2$ consists of all entire functions of exponential type less or equal to $\sigma$ whose restriction to the real line belong to the space $L^2(\mathbb{R}).$	The space $B_\sigma^2$ is a  reproducing kernel Hilbert space when endowed with the norm
%\begin{equation*}
%\left\| f\right\|^2 = \int^{\infty}_{-\infty}|f(t)|^2dt, \quad f\in B^2_{\sigma}.
%%\end{equation*}
%For each $w\in \mathbb{C},$ let $k_w$ denote the \textit{reproducing kernel} for $B_{\sigma}^2$ at %$w$; that is,
%%\begin{equation*}
%k_w(z)=\frac{\sin \sigma(z-\overline{w})}{\pi(z-\overline{w})}, \quad z\in \mathbb{C}.
%\end{equation*}
%This means that for each $w\in \mathbb{C},$ the evaluation functional $f\mapsto f(w)$ is bounded on $B^2_{\sigma}$ and $\left\langle f,k_w\right\rangle =f(w)$ for all $f\in B_{\sigma}^2.$ 

A bounded linear operator $T$ on a separable Hilbert space $\mathcal{H}$ is \emph{cyclic} with \emph{cyclic vector} $f\in\mathcal{H}$ if the linear span of the \emph{orbit} Orb$(f,T)=\{T^nf:n=0,1,2,\ldots\}$ is dense in $\mathcal{H}$. Similarly $T$ is \emph{supercyclic} if all scalar multiples of elements in Orb$(f,T)$ are dense in $\mathcal{H}$, and \emph{hypercyclic} if the orbit itself is dense. Cyclicity is a central theme in linear dynamics and has been studied widely due to its connection with the Invariant Subspace Problem (ISP). See the recent monographs \cite{Bayart-Matheron} and \cite{Linear Chaos} to learn more about linear dynamics. 

On the other hand a bounded linear operator $T$ on $\mathcal{H}$ is \emph{complex symmetric} if there exists an orthonormal basis for $\mathcal{H}$ with respect to which $T$ has a self-transpose matrix representation. An equivalent definition also exists. A \emph{conjugation} is a conjugate-linear operator $J:\mathcal{H}\to\mathcal{H}$ that satisfies the conditions \\

(a) $J$ is \emph{isometric}: $\langle Jf, Jg \rangle=\langle g,f \rangle$ $\forall$ $f,g\in\mathcal{H}$,\\

(b) $J$ is \emph{involutive}: $J^2=I$.\\ \\
We say that $T$ is $J$-\emph{symmetric} if $JT=T^*J$, and complex symmetric if there exists a conjugation $J$ with respect to which $T$ is $J$-symmetric. Complex symmetric operators are generalizations of complex symmetric matrices and of normal operators, and their study was initiated by Garcia, Putinar and Wogen \cite{Garcia 1, Garcia 2, Garcia 3, Garcia 4}.  If $JT=T^*J$ for some operator $T$ and conjugation $J$, then $T$ is cyclic (supercyclic) if and only if $T^*$ is cyclic (supercyclic). The conjugation $J$ serves as a bijection between the cyclic vectors of $T$ and $T^*$.



Let $\mathcal{S}$ be a  space of functions defined on a set $\Omega.$ 
A \emph{composition operator} $C_{\phi}$
on $\mathcal{S}$ with \emph{symbol} $\phi:\Omega\to\Omega$  is defined as
\[
C_{\phi}f=f\circ \phi, \quad f\in \mathcal{S}.
\]
The cyclicity of composition operators has been studied on several holomorphic function spaces (see \cite{Bayart, Gallardo, Noor}) and has become an active research topic within linear dynamics. In \cite{chacon2007composition}, Chácon, Chácon and Giménez initiated the study of composition operators on the classical Paley-Wiener space $B^2_{\pi}.$ They prove that $C_\phi$ is bounded on $B^2_{\pi}$ precisely when
\begin{align}\label{b1}
\phi(z)=az+b,\ \text{where} \ a\in \mathbb{R} \ \text{with} \ 0<|a|\leq 1 \ \text{and} \ b\in \mathbb{C}.
\end{align}
More recently Ikeda, Ishikawa and Yoshihiro \cite{Ikeda} show that this is true even in the general context of reproduction kernel Hilbert spaces of entire functions on $\mathbb{C}^n$. 

The cyclicity and complex symmetry of $C_\phi$ on the Hardy-Hilbert space $H^2(\mathbb{C}_+)$ of the half-plane were characterized by  Noor and Severiano \cite{Noor} for affine symbols $\phi$. In this article we study the cyclicity and complex symmetry of composition operators on the \textit{Paley-Wiener spaces} $B^2_{\sigma}$ for all $\sigma>0$. From results in \cite{Ikeda} it follows that the only bounded composition operators $C_\phi$ on $B^2_{\sigma}$ for $\sigma>0$ are those induced by symbols of the form \eqref{b1}. The Paley-Wiener spaces $B^2_\sigma$ are isometrically embedded into $H^2(\mathbb{C}_+)$ as so-called \emph{model subspaces} $K_\Theta$ (see \cite[page 305]{Nikolski}) defined as
\[
K_\Theta:=H^2(\mathbb{C}_+)\ominus\Theta H^2(\mathbb{C}_+) \ \ \mathrm{where} \ \ \Theta(z)=e^{i\sigma z}.
\]
Therefore one can view this work as the study of $C_\phi$ on some model subspaces of $H^2(\mathbb{C}_+)$. In contrast with $H^2(\mathbb{C}_+)$, our results show the existence of non-normal complex symmetric $C_\phi$ (see Theorem \ref{normal}), and that the cyclicity of $C_\phi$ and its adjoint $C_\phi^*$ in $B^2_\sigma$ depends on $\sigma>0$ (see Theorem \ref{cyclic}). In particular, we show that no $C_\phi$ is supercyclic on any $B^2_\sigma$. Finally we characterizese the reproducing kernels $(k_w)_{w\in\mathbb{C}}$ in $B^2_\sigma$ that are cyclic vectors for $C_\phi$ and show that this is closely related to the completeness of exponential sequences $(e^{i\lambda_n t})_{n\in\mathbb{N}}$ in $L^2[-\sigma,\sigma]$ (see Theorem \ref{cyclic vector}). The latter is a central question in non-harmonic Fourier analysis (see \cite{young}). Some of the main results are summarized in the following table.

\begin{table}[h]
\centering
\caption{Main results for $C_\phi$ on $B^2_\sigma$ where $\phi(z)=az+b$.}
\label{tab:example}
\begin{tabular}{|c|c|} % 'c' for centered columns; adjust as needed

\hline
$C_\phi$ cyclic & $a=1$ with $b\in\mathbb{C}\setminus\mathbb{R}$ or $0<|b|\leq\pi/\sigma$, $b\in\mathbb{R}$ \\
\hline
$C^*_\phi$ cyclic & $0<|a|<1$, $b\in\mathbb{C}$ or $C_\phi$ cyclic \\
\hline
$C_\phi$ supercyclic & Never \\
\hline
$C_\phi$ complex symmetric & $a=1$, $b\in\mathbb{C}$ or $a=-1$, $b\in\mathbb{C}$  \\ 
\hline
$C_\phi$ normal & $a=1$, $b\in\mathbb{C}$ or $a=-1$, $b\in\mathbb{R}$  \\

\hline
\end{tabular}
\end{table}
Before we begin, it is necessary to mention an error in \cite{chacon2007composition} during the computation of an adjoint formula for $C_\phi$ on $B^2_\pi$ which then leads to an incomplete description of normal $C_\phi$. This is discussed and corrected in Section 2.






%Let $\mathcal{S}$ be a  space of functions defined on a set $\Omega.$ 
%%A \emph{composition operator} $C_{\phi}$
%on $\mathcal{S}$ is an operator acting by composition to the right with a chosen self-map $\phi$ of $\Omega,$ i.e.,
%\[
%C_{\phi}f=f\circ \phi, \quad f\in \mathcal{S}.
%%\]
%The self-map $\phi:\Omega\to\Omega$ is called the \emph{symbol} of the composition operator $C_{\phi}.$ Operators of this type have been studied on
%%a variety of spaces, such as Hardy spaces, Fock spaces and Paley-Wiener spaces (see \cite{Alex, chacon2007composition, cowen2019composition, Ikeda, Le, Shapiro}). The main objective here is to relate the operator-theoretic properties of $C_{\phi}$ with the function-theoretic properties of the symbol $\phi.$ For instance, in \cite{chacon2007composition}, Chácon, Chácon and Giménez initiated the study of composition operators on the classical Paley-Wiener space $B^2_{\pi}.$ They prove that the only bounded composition operators $C_\phi$ on $B^2_{\pi}$ are those induced by affine symbols
%\[\phi(z)=az+b, \ \ a\in \mathbb{R},  \  0<|a|\leq 1,  \ b\in \mathbb{C}\] and that $B^2_{\pi}$ does not support compact composition operators. Recently Ikeda, Ishikawa and Yoshihiro \cite{Ikeda} show that both these results are true in the more general context of reproduction kernel Hilbert spaces of entire functions on $\mathbb{C}^n$. In contrast, composition operators on the Hardy space $H^2(\mathbb{D})$ may be compact and induced by non-affine symbols.


%In this article we focus on the complex symmetry of $C_\phi$ on the Paley-Wiener spaces $B^2_{\sigma}$ for $\sigma>0$. Let $\mathcal{L}(\mathcal{H})$ be the space of all bounded linear operators on a separable complex Hilbert space $\mathcal{H}.$ Recall that $T$ is \textit{normal} if $TT^*=T^*T,$ \textit{self-adjoint} if $T=T^*,$ \textit{unitary} if $TT^*=I=T^*T$ and \textit{complex symmetric} (or \emph{$C$-symmetric}) if there is a  conjugate-linear
%operator on $\mathcal{H}$ satisfying $C^2=I$ and $\left\langle Cx,Cy\right\rangle =\left\langle y,x\right\rangle $ for all $x,y\in \mathcal{H},$  for which $CT^*C=T.$ Such a conjugate-linear operator $C$ is called a \textit{conjugation}. It is well-known that the class of complex symmetric operators contains the normal operators.

 %Similar to $B^2_{\pi},$ the only bounded composition operators on $B^2_{\sigma}$ are those induced by affine symbols of the form 
%\begin{align}\label{b1}
%\phi(z)=az+b,\ \text{where} \ a\in \mathbb{R} \ \text{with} \ 0<|a|\leq 1 \ \text{and} \ b\in \mathbb{C}.
%\end{align}
 %Before we begin it is important to highlight a particular error in a result of \cite{chacon2007composition} that then leads to some other inaccuracies. More precisely in \cite[page 2209 ]{chacon2007composition} it is stated that if $\phi$ is as in \eqref{b1} with $0<a<1,$ then the adjoint of $C_\phi$ is given by
%\begin{align*}
%(C_{\phi}^*f)(z)=\frac{1}{a}f\left(\frac{z-\overline{b}}{a}\right), \quad f\in B^2_{\pi}.
%\end{align*}
%However, this formula fails since if $f$ has exponential type $\pi$ then $f_a(z)=f(z/a)$ has exponential type $\pi/a>\pi$ and hence $C_\phi^*$ is not well-defined as an operator on $B^2_{\pi}$. As a result, an erroneous characterization of normal $C_{\phi}$ on $B^2_{\pi}$ is obtained (see \cite[Prop. 2.6]{chacon2007composition}).



%\begin{table}[h]
%\centering
%\caption{Main results where $\phi(z)=az+b$.}
%\label{tab:example}
%\begin{tabular}{|c|c|} % 'c' for centered columns; adjust as needed
%\hline
%$C_\phi$ complex symmetric & $a=1$, $b\in\mathbb{C}$ or $a=-1$, $b\in\mathbb{C}$  \\ 
%\hline
%$C_\phi$ normal & $a=1$, $b\in\mathbb{C}$ or $a=-1$, $b\in\mathbb{R}$  \\
%\hline
%$C_\phi$ self-adjoint & $a=1$, $b\in \mathrm{i}\mathbb{R}$ or $a=-1$, $b\in \mathbb{R}$ \\ 
%\hline
%$C_\phi$ unitary & $a=1$, $b\in\mathbb{R}$ or $a=-1$, $b\in\mathbb{R}$ \\
%\hline
%\end{tabular}
%\end{table}






%====================================================================================================
\section{Preliminaries}
%====================================================================================================
An entire function $f$ is of \textit{exponential type} if the inequality $
|f(z)|\leq Ae^{B|z|}$ holds for all $z\in \mathbb{C}$
and for some constants $A,B>0$. The exponential type $\sigma$ of $f$ is defined as the infimum of all $B>0$ for which this inequality holds and can be determined by 
\begin{align*}
\sigma=\limsup\limits_{r\rightarrow\infty}\frac{\log M_f(r)}{r} \ \  \mathrm{where} \ \ M_f(r)=\max_{|z|=r}|f(z).
\end{align*}


\subsection{Paley-Wiener spaces \texorpdfstring{$B^2_{\sigma}$}{lg}}  For $\sigma>0,$ the \emph{Paley-Wiener space} $B_\sigma^2$ consists of all entire functions of exponential type $\leq\sigma$ whose restrictions to $\mathbb{R}$ belong to $L^2(\mathbb{R}).$	The space $B_\sigma^2$ is a  reproducing kernel Hilbert space when endowed with the norm
\begin{equation*}
\left\| f\right\|^2 = \int^{\infty}_{-\infty}|f(t)|^2\mathrm{d}t, \quad f\in B^2_{\sigma}.
\end{equation*}
For each $w\in \mathbb{C},$ let $k_w$ denote the reproducing kernel for $B_{\sigma}^2$ at $w$ defined by
\begin{equation*}
k_w(z)=\frac{\sin \sigma(z-\overline{w})}{\pi(z-\overline{w})}, \quad z\in \mathbb{C}
\end{equation*}
and which satisfies the basic relation $\left\langle f,k_w\right\rangle =f(w)$ for all $f\in B_{\sigma}^2.$ The Paley-Wiener Theorem states that any $f\in B^2_\sigma$ can be represented as the \emph{inverse} Fourier transform 
\[
f(z)=(\mathfrak{F}^{-1}F)(z):=\frac{1}{\sqrt{2\pi}}\int_{-\sigma}^\sigma F(t)e^{izt}dt
\]
of some $F\in L^2[-\sigma,\sigma]$ where the Fourier transform $\mathfrak{F}:B^2_\sigma\to L^2[-\sigma,\sigma]$ is an isometric isomorphism. We denote $\widehat{f}:=\mathfrak{F}f$ for $f\in B^2_\sigma$ and in particular $(\widehat{k_w})(t)=e^{-i\bar{w}t}$ for $w\in\mathbb{C}$.
The monographs \cite{Branges} and \cite{young} may be consulted for more information about these spaces.
%===================================================================================================
\subsection{Composition operators}
%===================================================================================================
The composition operator $C_{\phi}$ is bounded on $B^2_{\sigma}$ if and only if  its inducing symbol $\phi$ has the form 
\begin{align}\label{symbol}
\phi(z)=az+b,\ \text{where} \ a\in \mathbb{R} \ \text{with} \ 0<|a|\leq 1 \ \text{and} \ b\in \mathbb{C}
\end{align}	
%(the proof of this fact is similar to what was done in \cite[Theorem 2.4.]{chacon2007composition})
with the usual action of $C_{\phi}^*$ on reproducing kernels determined by $C_{\phi}^*k_w=k_{\phi(w)}$ for $w\in \mathbb{C}.$ For each $n\in \mathbb{N},$ the $n$-iterate of the self-map $\phi:\mathbb{C}\rightarrow \mathbb{C}$ is denoted by $\phi^{[n]}$ where
\begin{align}\label{iterate}
\phi^{[n]}(z)= \begin{cases}z+n b, & \text{if} \ a=1, \\ 
a^{n} z+\frac{\left(1-a^{n}\right)}{1-a}b,  & \text{if} \ a \neq 1 
\end{cases}
\end{align}	
and $C_{\phi}^n=\C_{\phi^{[n]}}.$ We see from \eqref{iterate} that if $0<|a|<1$ and $\alpha:=\frac{b}{1-a},$ then $\alpha$ is an attractive fixed point of $\phi$, that is,  $\phi^{[n]}(z)\rightarrow \alpha$ as $n\rightarrow \infty$ for all $z\in \mathbb{C}.$ 








\subsection{Adjoints of Composition operators} In \cite{chacon2007composition} the authors show that $C_\phi$ and $C_\phi^*$ on $B^2_\pi$ are unitarily equivalent, via the Fourier transform $\mathfrak{F}$, to a pair of weighted composition operators $\widehat{C}_{\phi}$ and $\widehat{C}_{\phi}^*$ on $L^2[-\pi,\pi]$ respectively. For simplicity they assume $a>0$ which is sufficient for their results. However, as we shall see in the next section, the cases $a>0$ and $a<0$ lead to very distinct outcomes for the complex symmetry and cyclicity of $C_\phi$. Therefore for the sake of completeness we provide a proof of their result for all $0<|a|\leq 1$.




\begin{prop}\label{adjoint formula} If $\phi(z)=az+b$ where $a\in \mathbb{R}$ with $0<|a|\leq 1$, $b\in \mathbb{C},$ then $C_{\phi}$ on $ B^2_{\sigma}$ is unitarily equivalent to a weighted composition operator $\widehat{C}_{\phi}$ on $L^2[-\sigma,\sigma]$ defined by 
\[(\widehat{C}_{\phi}F)(t)=\frac{1}{|a|}\chi_{(-|a|\sigma, |a|\sigma)}(t)e^{\frac{\mathrm{i}bt}{a}}F\left(\frac{t}{a}\right)\] 
where $\chi_{(c,d)}$ denotes a characteristic function. Moreover we have $(\widehat{C}_{\phi}^*F)(t)=\overline{e^{ibt}}F(at).$
\end{prop}
\begin{proof} We first assume $0<a\leq 1.$ For each $f\in B^2_{\sigma},$ we have
\begin{align}
(C_{\phi}f)(z)&=f(az+b)=\frac{1}{\sqrt{2\pi}}\int^\sigma_{-\sigma}\widehat{f}(t)e ^{i(az+b)t}dt\nonumber\\
&=\frac{1}{\sqrt{2\pi}}\int^\sigma_{-\sigma}\widehat{f}(t)e ^{iazt}e^{ibt}dt=\frac{1}{\sqrt{2\pi}}\int^{a\sigma}_{-a\sigma}\frac{1}{a}\widehat{f}\left(\frac{s}{a}\right)e ^{isz}e^{ib\left(\frac{s}{a}\right)}ds\nonumber\\
&=\frac{1}{\sqrt{2\pi}}\int^{\sigma}_{-\sigma}\frac{1}{a}\chi_{(-a\sigma, a\sigma)}(t)e^{\frac{ibs}{a}}\widehat{f}\left(\frac{s}{a}\right)e ^{isz}ds\nonumber\\
&=\frac{1}{\sqrt{2\pi}}\int^{\sigma}_{-\sigma}(\widehat{C}_{\phi}\widehat{f})(s)e ^{isz}ds=(\mathfrak{F}^{-1}\widehat{C}_\phi\mathfrak{F}f)(z)\nonumber
\end{align}
which gives $\mathfrak{F}C_{\phi}= \widehat{C}_{\phi}\mathfrak{F}$. For $-1\leq a <0$, first consider the symbol $\eta(z)=-z$. Then the simple change of variables $s=-t$ again gives
\begin{align}\label{eqc3}
(C_{\eta}f)(z)&=f(-z)=\frac{1}{\sqrt{2\pi}}\int^{\sigma}_{-\sigma}(\widehat{C}_{\eta}\widehat{f})(s)e ^{isz}ds=(\mathfrak{F}^{-1}\widehat{C}_\eta\mathfrak{F}f)(z)
\end{align}
and hence $\mathfrak{F}C_{\eta}= \widehat{C}_{\eta}\mathfrak{F}.$  Now let $\psi(z)=-az+b$ and note that $C_{\phi}=C_{\eta}C_{\psi}.$ Therefore 
\begin{align*}
\widehat{C}_{\eta}\widehat{C}_{\psi}=(\mathfrak{F}C_{\eta}\mathfrak{F}^{-1})(\mathfrak{F}C_{\psi}\mathfrak{F}^{-1})=\mathfrak{F}C_{\phi}\mathfrak{F}^{-1}
\end{align*}
and to verify that indeed $\widehat{C}_\phi=\widehat{C}_{\eta}\widehat{C}_{\psi}$, one easily sees that
\begin{align*}
(\widehat{C}_{\eta}\widehat{C}_{\psi}F)(t)&=\frac{1}{-a}e^{-\frac{ibt}{-a}}\chi_{(a\sigma, -a\sigma)}(-t)F\left(\frac{-t}{-a}\right)
=(\widehat{C}_{\phi}F)(t)
\end{align*}
for all $F\in L^2[-\sigma,\sigma]$. Therefore $\mathfrak{F}C_{\phi}= \widehat{C}_{\phi}\mathfrak{F}$ for all $-1\leq a<0$ as well. Since the Fourier transform $\mathfrak{F}:B^2_\sigma\to L^2[-\sigma,\sigma]$ is an isometric isomorphism, the proof is complete. The adjoint 
formula follows by the change of variables $t=s/a$ as follows:
\[
\langle \widehat{C}_{\phi}^*F, G\rangle=\int^\sigma_{-\sigma}\overline{e^{ibt}}F(at)\overline{G(t)}dt=\int^{|a|\sigma}_{-|a|\sigma}F(s)\overline{\frac{1}{|a|}e^{\frac{\mathrm{i}bs}{a}}G\left(\frac{s}{a}\right)}ds=\langle F, \widehat{C}_{\phi}G\rangle
\]
for all $F,G\in L^2[-\sigma,\sigma]$. This completes the proof of the result.
\end{proof}
In \cite[p. 2209]{chacon2007composition} it is claimed that if $\phi=az+b$ with $0<a<1,$ then the adjoint of $C_\phi$ is 
\begin{align*}
(C_{\phi}^*f)(z)=\frac{1}{a}f\left(\frac{z-\overline{b}}{a}\right), \quad f\in B^2_{\pi}.
\end{align*}
However this cannot be correct since if $f$ has exponential type $\pi$ then $f(z/a)$ must have type $\pi/a>\pi$. Hence $C_\phi^*$ is not well-defined on $B^2_{\pi}$. This unfortunately leads to an incomplete description of normal composition operators (compare \cite[Prop. 2.6]{chacon2007composition} with Table $1$).

%===================================================================================================
%\section{The norm of \texorpdfstring{$C_{\phi}$}{lg}}\label{s2}
%==================================================================================================
  

%Now we move on to the work of determining the norm of each bounded composition operator and its spectral radius on $B^2_{\sigma}.$ Here the main idea is to write $\phi(z)=az+b$ where $a\in \mathbb{R}$ with $0<|a|\leq 1$ as follows  
%\begin{align*}
%\phi=\psi\circ \varphi
%\end{align*}
%where $\varphi(z)= az$ and $\psi(z)=\pm z+b,$ and then %we estimate the norms of the auxiliary bounded composition operators $C_{\varphi}$ and $C_{\psi}.$

%\begin{rem}\label{r1}
%In \cite[page 2209]{chacon2007composition} is stated that if $\phi(z)=az+b$ where $a\in \mathbb{R}$ with $0<|a|\leq 1$ and $b\in \mathbb{C},$ then the adjoint of $C_{\phi}$ on $B^2_{\pi}$ is $(C_{\phi}^*f)(z)=\frac{1}{a}f\left( \frac{z-\overline{b}}{a}\right).$ However this statement is false, because if $f$ has exponential type $\pi$ and $|a|<1$ then $g$ defined by $g(z)=f(z/a)$ has exponential type $\pi/|a|>\pi.$
%\end{rem}
 
%\begin{lem}\label{rotation} Let $\phi(z)=az$ where $a\in \mathbb{R}$ with $0<|a|\leq 1.$ Then 
%$\|C_{\phi} f\|=|a|^{-1/2}\|f\|$ for each $f\in B^2_{\sigma}.$   \end{lem}
 
% \begin{proof} We first assume $0<a\leq 1.$ Then for each $f\in B^2_{\sigma},$ we have 
%\begin{align*}
%\|C_{\phi}f\|^2=\int^{\infty}_{-\infty}|f(at)|^2dt=\frac{1}{a}\int^{\infty}_{-\infty}|f(t)|^2dt=\frac{1}{a}\|f\|^2,
%\end{align*}
%which implies $\|C_{\phi}f\|^2=a^{-1/2}\|f\|.$ For $-1\leq a<0,$ we consider the symbols $\psi(z)=-z$ and $\varphi(z)=-az.$ By Lemma \ref{isometry} the operator $C_{\psi}$ is an isometric isomorphism, and from what we have shown so far $\|C_{\varphi}f\|=(-a)^{-1/2}\|f\|$ for each $f\in B^2_{\sigma}.$ Since $C_{\phi}=C_{\psi}C_{\varphi},$ it follows that
%\begin{align*}
%%\|C_{\phi}f\|=\|C_{\psi}C_{\varphi}f\|=\|C_{\varphi}f\|=(-a)^{-1/2}\|f\|
%%\end{align*}
%for each $f\in B^2_{\sigma}.$
 %\end{proof}


 
%\begin{thm}\label{norm} Let $\phi(z)=az+b$ where $a\in \mathbb{R}$ with $0<|a|\leq 1.$ Then $\|C_{\phi}f\|=|a|^{-1/2}\|f\|$ for each $f\in B^2_{\sigma}.$ In particular, $\|C_{\phi}\|=|a|^{-1/2}.$
%\end{thm}
%\begin{proof}  We first consider $0<a\leq 1,$ in this case let  $\psi(z)=z+b$ and $\varphi(z)=az$ be the auxiliary symbols. Since $\phi=\psi\circ \varphi$ and $C_{\psi}$ is an isometry, we get 
%\begin{align}
%\|C_{\phi}f\|=\|C_{\varphi}C_{\psi}f\|=a^{-1/2}\|C_{\psi}f\|=a^{-1/2}\|f\|,
%\end{align}
%for all $f\in B^2_{\sigma}.$ Now we consider $-1\leq a<0.$ Let $\varphi(z)=-az$ and $\psi(z)=-z-b/a$ be the auxiliary symbols. Since $\phi=\varphi\circ\psi $ and $C_{\psi}$ is an isometric isomorphism, we get $C_{\psi}^{-1}C_{\phi}=C_{\varphi}.$  From what we have shown so far $\|C_{\varphi}f\|=(-a)^{1/2}\|f\|$ for all $f\in B^2_{\sigma}$ and hence 
%\begin{align*}
%\|C_{\phi}f\|=\|C_{\psi}^{-1}C_{\phi}f\|=\|C_{\varphi}f\|=(-a)^{-1/2}\|f\|.
%\end{align*}
%This completes the proof.
%\end{proof}


%\begin{cor} If $C_{\phi}$ is a bounded composition operator on $B^2_{\sigma},$ then $C_{\phi}$ has closed range.
%\end{cor}

%\begin{cor} Let $\phi(z)=az+b$ where $a\in \mathbb{R}$ with $0<|a|\leq 1.$ Then $C_{\phi}$ is an isometry on $B^2_{\sigma}$ if and only if $|a|=1.$
%\end{cor}


%\begin{cor} If $C_{\phi}$ is a bounded composition operator on $B^2_{\sigma},$ then $r(C_{\phi})=\|C_{\phi}\|.$
%\end{cor}
%\begin{proof} Let $C_{\phi}$ be a bounded composition operator on $B^2_{\sigma},$ then $\phi(z)=az+b$ where $a\in \mathbb{R}$ with $0<|a|\leq1.$ By \eqref{iterate} and Theorem \ref{norm}, we obtain $\|C_{\phi}^n\|=|a|^{-n/2}$ for each $n\in \mathbb{N}.$  So taking $n^{th}$ roots and letting $n$ go to infinity we get $r(C_{\phi})=|a|^{-1/2}=\|C_{\phi}\|.$
%\end{proof}

%\textcolor{red}{In \cite[Section 2]{Barbara}, Cowen and Barbara proposed the following question for bounded composition operators on $H^2(\mathbb{D}):$ when the norm of $\|C_{\phi}\|$ can recover solely from the action of $C_{\phi}$ or $C_{\phi}^*$ on the reproducing kernel functions $k_z$ for $H^2(\mathbb{D});$ that is,
%is either
%\begin{align}\label{supremum}
%\|C_{\phi}\|=\sup_{z\in \mathbb{D}} \frac{\|C_{\phi}^*k_z\|}{\|k_z\|} \quad \text{or} \quad \|C_{\phi}\|=\sup_{z\in \mathbb{D}} \frac{\|C_{\phi}k_z\|}{\|k_z\|}
%\end{align}
%is true? Appel, Bourdon, and Thrall \cite{Appel} and Bourdon and Retsek \cite{Bourdon} show that the answer to both of these questions is no in general. The version of this question for composition operators on the Hardy space of right half-plane $H^2(\mathbb{C}_+)$ was answered positively by Elliot and Jury \cite[Corollary 3.2.]{Elliot}. We finish this section by studying the version of the problem for composition operators on the Paley-Wiener space $B^2_{\sigma}.$ }



%\begin{lem}\label{norm of kernel}  For $w\in\C,$ the norm $\|k_w\|$ is given by
%\begin{align*}
%\|k_w\|^2=\begin{cases}\displaystyle\frac{\operatorname{sinh}2\sigma\mathrm{Im}(w)}{2\pi\mathrm{Im}(w)}, & \text{if}\ \mathrm{Im}(w)\neq 0, \\ 
%\displaystyle \frac{\sigma}{\pi}, & \text{if}\  \mathrm{Im}(w) = 0.
%\end{cases}
%\end{align*}
%\end{lem}
%\begin{proof}
%Let $w\in \mathbb{C}.$ If $\mathrm{Im}(w)\neq 0,$ we have
%\begin{align*}
%\|k_w\|^2= k_{w}(w)=\frac{\operatorname{sin}\sigma(w-\overline{w})}{\pi(w-\overline{w})}=\frac{\operatorname{sin}2\mathrm{i}\sigma\mathrm{Im}(w)}{2\mathrm{i}\pi\mathrm{Im}(w)}=\frac{\mathrm{i}\operatorname{sinh}2\sigma\mathrm{Im}(w)}{2\mathrm{i}\pi \mathrm{Im}(w)}=\frac{\sinh 2\sigma \im(w)}{2\pi\mathrm{Im}(w)},
%\end{align*}
%while $\mathrm{Im}(w)=0$ gives $\|k_w\|^2=\sigma/\pi.$ 
%\end{proof}


%\begin{prop}\label{kernel supremum} If $C_{\phi}$ is a bounded composition operator on $B^2_{\sigma}$ then 
%\begin{align*}
%\|C_{\phi}\|=\sup_{w\in \mathbb{C}}\frac{\|C_{\phi}^*k_w\|}{\|k_w\|}.
%\end{align*}
%\end{prop}
%\begin{proof}  Let $C_{\phi}$ be a bounded composition operator on $B^2_{\sigma},$ then $\phi(z)=az+b$ where $a\in \mathbb{R}$ with $0<|a|\leq1.$ For the sake of simplicity we assume $0 < a \leq 1.$ Let $b=0.$ For $w\notin \mathbb{R},$ Lemma \ref{norm of kernel} gives
%\begin{align}
%\frac{1}{a}=\|C_{\phi}\|^2\geq \frac{\|C_{\phi}^*k_w\|^2}{\|k_w\|^2}=\frac{\|k_{\phi(w)}\|^2}{\|k_w\|^2}=\frac{1}{a}\frac{\sinh 2\sigma a\mathrm{Im}(w)}{\sinh 2\sigma\mathrm{Im}(w)}.
%\end{align}
%Since sinh$(x)$ is non-decreasing on the entire real line, we obtain
%\begin{align*}
%\|C_{\phi}\|^2&\geq \sup \left\{\frac{\|C_{\phi}^*k_w\|^2}{\|k_w\|^2}:w\in \mathbb{C}\right\}
%=\frac{1}{a}\sup \left\{ \frac{\sinh 2\sigma a\mathrm{Im}(w)}{\sinh 2\sigma\mathrm{Im}(w)}:w \in \mathbb{C} \right\}\\
%& \geq \frac{1}{a}\sup \left\{ \frac{\sinh 2\sigma a\mathrm{Im}(w)}{\sinh 2\sigma\mathrm{Im}(w)}:\mathrm{Im}(w)<0\right\}
%\geq \frac{1}{a}.
%\end{align*}
%For the general case, we consider the auxiliary symbols $\varphi(z)=az$ and $\psi(z)=z-b.$ Then $C_{\psi}$ is an isometric isomorphism and $C_{\phi}C_{\psi}=C_{\varphi}.$ 
%From what we have shown so far, we have
%\begin{align*}
%\|C_{\phi}\|=\|C_{\varphi}C_{\psi}^{-1}\|=\|C_{\varphi}\|=\sup_{w\in \mathbb{C}}\frac{\|C_{\varphi}^*k_w\|}{\|k_w\|}=\sup_{w\in \mathbb{C}}\frac{\|C_{\psi}^*C_{\phi}^*k_w\|}{\|k_w\|}=\sup_{w\in \mathbb{C}}\frac{\|C_{\phi}^*k_w\|}{\|k_w\|}
%\end{align*}
%as desired.
%\end{proof}
%In this section, we use the reproducing kernels to obtain a lower estimate for the essential norm for bounded composition operators, and then we show that no bounded composition operator on $B^2_{\sigma}$ is compact. 

%\begin{lem} The sequence $(k_{\frac{n\pi}{\sigma}})_{n\in \mathbb{N}}$ is orthogonal on $B^2_{\sigma}.$
%\end{lem}
%\begin{proof} For $n, m\in \mathbb{N}$ with $n\neq m,$ we have 
%\begin{align*}
%\langle k_{\frac{n\pi}{\sigma}}, k_{\frac{m\pi}{\sigma}}\rangle=k_{\frac{n\pi}{\sigma}}\left(\frac{m\pi}{\sigma}\right)=
%\frac{\sin \sigma \left( \frac{m\pi}{\sigma}- \frac{n\pi}{\sigma}\right)}{\pi \left( \frac{m\pi}{\sigma}- \frac{n\pi}{\sigma}\right)}=\frac{\sin  \left(m-n \right)\pi}{\pi \left( \frac{m\pi}{\sigma}- \frac{n\pi}{\sigma}\right)}=0
%\end{align*}
%as desired.
%\end{proof}



%\begin{prop} If $T:B^2_{\sigma}\rightarrow %\mathcal{H}$ is a bounded operator which satisfies
%\begin{align*}
%\|T\|=\sup_{w\in \mathbb{C}}\frac{\|T^*k_w\|}{\|k_w\|} \quad \text{or} \quad \|T\|=\sup_{w\in \mathbb{C}}\frac{\|Tk_w\|}{\|k_w\|} 
%\end{align*}
%then $\|T\|=\|T\|_e.$
%\end{prop}
%\begin{proof} The inequality $\|T\|_{e}\leq \|T\|$ is immediate. For the reverse inequality, 
 %Let $\epsilon>0,$ then there exists a compact operator $T_{\epsilon}$ such that $\epsilon+\|T\|_{e}>\|T-T_{\epsilon}\|.$ Since $\|T-T_{\epsilon}\|=\|T^*-T_{\epsilon}^*\|$ and $T_{\epsilon}^*$ is compact, for simplicity we can assume $\|T\|=\displaystyle \sup_{w\in \mathbb{C}}\frac{\|Tk_w\|}{\|k_w\|} $ and hence
%\begin{align}\label{compact}
%\epsilon+\|T\|_e&
%\geq \sup_{w\in \mathbb{C}}\frac{\|Tk_w\|}{\|k_w\|}-\frac{\|T_{\epsilon}k_w\|}{\|k_w\|}=\|T\|-\frac{\|T_{\epsilon}k_w\|}{\|k_w\|}
%\end{align}
%for all $w\in \mathbb{C}.$  Since $k_{\frac{n\pi}{\sigma}}/\|k_{\frac{n\pi}{\sigma}}\|$ tend weakly to zero as $n\rightarrow \infty$ and $T_{\epsilon}$ is compact, it follows from \eqref{compact} that $\epsilon+\|T\|_e
%\geq \|T\|.$ This completes the proof.
%\end{proof}


%\begin{proof} The inequality $\|C_{\phi}\|_{e}\leq \|C_{\phi}\|$ is immediate. For the reverse inequality, let $\epsilon>0,$ then there exists a compact operator $T_{\epsilon}$ such that $\epsilon+\|C_{\phi}\|_{e}>\|C_{\phi}-T_{\epsilon}\|\geq \|C_{\phi}f\|-\|T_{\epsilon}f\|$ for all $f\in B^2_{\sigma}$ with $\|f\|=1.$ Hence, Proposition \ref{kernel supremum} we get
%\begin{align}\label{compact}
%\epsilon+\|C_{\phi}\|_e&
%\geq \sup_{w\in \mathbb{C}}\frac{\|C_{\phi}^*k_w\|}{\|k_w\|}-\frac{\|T_{\epsilon}k_w\|}{\|k_w\|}=\|C_{\phi}\|-\frac{\|T_{\epsilon}k_w\|}{\|k_w\|}
%\end{align}
%for each $w\in \mathbb{C}.$  Since $k_{\frac{n\pi}{\sigma}}/\|k_{\frac{n\pi}{\sigma}}\|$ tend weakly to $0$ as $n\rightarrow \infty$ and $T_{\epsilon}$ is compact, it follows from \eqref{compact} that $\epsilon+\|C_{\phi}\|_e
%\geq \|C_{\phi}\|.$ This completes the proof.
%\end{proof}


 %The next result shows that bounded composition operators on $B^2_{\sigma}$ have a similar behavior to the composition operators on $H^2(\mathbb{C}_+)$, in the sense that its norm coincides with itself essential norm. 


%\begin{cor} If $C_{\phi}$ is a bounded composition operator on $B^2_{\sigma}$ then $ \|C_{\phi}\|_e =\|C_{\phi}\|.$
%\end{cor}



%\begin{cor}
%There are no compact composition operators on $B^2_{\sigma}.$
%\end{cor}

%The fact that $B^2_{\sigma}$ does not support compact composition operators can be seen as a particular case of a more general context, as shown in \cite[Proposition 3.2]{Ikeda}. 

%================================================================================================
\section{Cyclicity and complex symmetry of  \texorpdfstring{$C_{\phi}$}{lg}}\label{s6}
%==================================================================================================
One of the main leitmotifs of this article is to show how results about cyclicity can lead to results about complex symmetry and vice versa. This hinges on a simple observation. If $TJ=JT^*$ for some operator $T$ and conjugation $J$, then $T$ is cyclic if and only if $T^*$ is cyclic. The conjugation $J$ serves as a bijection between the cyclic vectors of $T$ and $T^*$. We demonstrate this by first dealing with the case when $0<|a|<1$.

\begin{prop}\label{case a<1} Let $\phi(z)=az+b$ where $a\in \mathbb{R}$, $0<|a|<1$ and $b\in\mathbb{C}$. Then $C_\phi$ is not cyclic whereas $C_\phi^*$ is cyclic on $B^2_\sigma$ for $\sigma>0$. Hence $C_\phi $ is not complex symmetric.
\end{prop}
\begin{proof} We know that $C_\phi$ on $B^2_\sigma$ is unitarily equivalent to $\widehat{C}_\phi$ on $L^2[-\sigma,\sigma]$ by Proposition \ref{adjoint formula}. So we show that $\widehat{C}_\phi$ is not cyclic. For any $F\in L^2[-\sigma,\sigma]$ we have $\widehat{C}_\phi F=\chi_{(-|a|\sigma, |a|\sigma)}G$ for some $G\in L^2[-\sigma,\sigma]$. So every element in $\mathrm{span}\{\widehat{C}_\phi^n F:n=0,1\ldots\}$ must have the form $cF+H$ where $H$ vanishes ouside $(-|a|\sigma, |a|\sigma)$. In other words, the closure of this span consists only of functions that coincide with a constant multiple of $F$ outside $(-|a|\sigma, |a|\sigma)$. This clearly makes the cyclicity of $\widehat{C}_\phi$ and hence of $C_\phi$ impossible. For $C^*_\phi$ we show that every reproducing kernel $k_w$ is a cyclic vector. This is because $(C^{*}_\phi )^nk_w=C^*_{\phi^{[n]}}k_w=k_{\phi^{[n]}(w)}$ and $\phi^{[n]}(w)$ converges to the attractive fixed point $\alpha:=\frac{b}{1-a}$ of $\phi$. So any $f\in B^2_\sigma$ orthogonal to the orbit of $k_w$ under $C^*_\phi$ must vanish on a sequence with a limit point and hence $f\equiv 0$. Therefore the span of the orbit of every $k_w$ must be dense in $B^2_\sigma$ and hence $C_\phi^*$ is cyclic. It then follows that $C_\phi$ is not complex symmetric by the discussion above. 
\end{proof}

By Proposition \ref{case a<1} we need only consider the cases $a=\pm1$ for the following result. Note that in both cases ($J_af)(z)=\overline{f(-a\bar{z})}$ defines a conjugation on $B^2_\sigma$ for $\sigma>0$.








\begin{prop}\label{normal}Let $\phi(z)=az+b$ where $a=\pm1$ and $b\in \mathbb{C}$. Then $C_{\phi}$ on $B^2_{\sigma}$ is
\begin{enumerate}
\item always complex symmetric,
\item normal if and only if $a=1$, $b\in\mathbb{C}$ or $a=-1, b\in \mathbb{R}$,
\item self-adjoint if and only if $a=1,b\in i\mathbb{R}$ or $a=-1,b\in \mathbb{R},$
\item unitary if and only if $b\in \mathbb{R}.$
\end{enumerate}
Moreover $C_\phi$ is $J_a$-symmetric on $B^2_{\sigma}$.
\end{prop}
\begin{proof} We first note that when $a=-1$ we have $C^2_\phi=I$ and hence $C_\phi$ is complex symmetric since every operator that is algebraic of order $2$ is complex symmetric by \cite[Thm. 2]{Garcia 1}. In general by Proposition \ref{adjoint formula} we have \[
(\widehat{C}_\phi F)(t) =e^{\frac{ibt}{a}}F(t/a) \ \ \ \mathrm{and} \ \ \ (\widehat{C}^*_\phi F)(t) =e^{\overline{ibt}}F(at)
\]
for $F\in L^2[-\sigma,\sigma]$. So when $a=1$ we see that $\widehat{C}_\phi=M_{e^{ibt}}$ is a multiplication operator which is normal. This shows that $C_\phi$ is always complex symmetric and gives $(1)$. If $a=-1$, then $\widehat{C}_\phi \widehat{C}_\phi^*=\widehat{C}_\phi ^*\widehat{C}_\phi$ precisely when $e^{-ibt}e^{-\overline{ibt}}=e^{\overline{ibt}}e^{ibt}$ or when $e^{-i2\mathrm{Im}(b)t}=e^{i2\mathrm{Im}(b)t}$, and this holds only when $b\in\mathbb{R}$. This gives $(2)$ and one easily obtains $(3)$ and $(4)$ similarly. Finally we show that $C_\phi$ is $J_a$-symmetric by first noting that
\[
(J_a C_\phi J_af)(z)=\overline{(C_\phi J_af)(-a\bar{z})}=\overline{(J_af)(-\bar{z}-a\bar{b})}=f(az-a\bar{b})=(C_\psi f)(z)
\]
where $\psi(z)=az-a\bar{b}$. We claim that $C_\psi=C_\phi^*$ and this follows by Proposition \ref{adjoint formula} because \[(\widehat{C}_{\psi}F)(t)=e^{-i\bar{b}t}F(t/a)=\overline{e^{ibt}}F(at)=(\widehat{C}_{\phi}^*F)(t)\] for all $F\in L^2[-\sigma,\sigma]$. Therefore $J_a C_\phi J_a=C_\phi^*$ and completes the proof of the result.
\end{proof}



Our main result characterizes the cyclicity of $C_\phi$ and $C_\phi^*$ simultaneously using complex symmetry, and then shows that no $C_\phi$ is supercyclic on $B^2_\sigma$ for all $\sigma>0$ via normality. The latter was proved for the case $\sigma=\pi$ in \cite[Thm. 2.7]{chacon2007composition} using a lengthier argument.

\begin{thm}\label{cyclic}Let $\phi(z)=az+b$ where $a\in\mathbb{R}$, $0<|a|\leq 1$ and $b\in \mathbb{C}$. Then $C_{\phi}$ on $B^2_\sigma$ is
\begin{enumerate}
\item cyclic if and only if $a=1$ with $b\in\mathbb{C}\setminus\mathbb{R}$ or $0<|b|\leq\pi/\sigma$, $b\in\mathbb{R}$, and
\item never supercyclic.
\end{enumerate}
Moreover $C_\phi^*$ is cyclic if and only if $0<|a|<1$ or when $C_\phi$ is cyclic.
\end{thm}
\begin{proof} The case $0<|a|<1$ has been dealt with in Proposition \ref{case a<1} showing that $C_\phi$ is not cyclic or supercyclic. So first note that when $a=-1$ then $C_\phi^2f=f$ and hence $C_\phi$ cannot be cyclic or supercyclic as the orbit of any $f\in B^2_\sigma$ contains at most two elements $f$ and $f\circ\phi$. So let $a=1$. Then $C_\phi$ is normal and normal operators are never supercyclic \cite[p. 564]{Hilden}. Therefore this proves $(2)$. For cyclicity we first observe that $\widehat{C}_\phi=M_{e^{ibt}}$ is a multiplication operator on $L^2[-\sigma,\sigma]$. A consequence of the spectral theory for normal operators is that a multiplication operator $M_\Phi$ on $L^2(\mu)$ (where $\Phi\in L^\infty(\mu)$ and $\mu$ is a compactly supported measure on $\mathbb{C}$) is cyclic if and only if $\Phi$ is injective on a set of full measure (see \cite[Thm. 1.1 and Prop. 1.3]{Ross}). Therefore we need to determine when $\Phi(t)=e^{ibt}$ is injective almost everywhere on $[-\sigma,\sigma]$. For $b\in\mathbb{C}\setminus\mathbb{R}$ we see that $|\Phi(t)|=e^{-\mathrm{Im}(b)t}$ which is clearly injective on all of $[-\sigma,\sigma]$ and hence so is $\Phi$. For the case when $b\in\mathbb{R}$ note that $\Phi$ is $2\pi/|b|$ periodic. It follows that $\Phi$ is injective on $(-\sigma,\sigma)$ (hence a.e on $[-\sigma,\sigma]$) precisely when the period is greater than or equal to the length of the interval, that is, precisely when
\[
\frac{2\pi}{|b|}\geq 2\sigma \ \ \ \mathrm{equivalently} \ \ \ 0<|b|\leq\frac{\pi}{\sigma}.
\]
The assertion about cyclicity of $C_\phi^*$ follows by Proposition \ref{case a<1} and complex symmetry. 
\end{proof}
Finally we ask the following question: \emph{Are any of the reproducing kernels $(k_w)_{w\in\mathbb{C}}$ cyclic vectors for $C_\phi$}? The answer reveals an interesting dichotomy and a connection with the question of completeness of exponential sequences in $L^2[-\sigma,\sigma]$. 

\begin{thm}\label{cyclic vector}Let $\phi(z)=z+b$ where $b\in \mathbb{C}$. Then the kernels $(k_w)_{w\in\mathbb{C}}$ in $B^2_\sigma$ are 
\begin{enumerate}
\item all cyclic vectors for $C_\phi$ when $b\in\mathbb{C}\setminus\mathbb{R}$ or $0<|b|<\pi/\sigma$, $b\in\mathbb{R}$, and
\item never cyclic for $C_\phi$ when $|b|=\pi/\sigma$. 
\end{enumerate}
\end{thm}
\begin{proof} We first consider the case when $b\in\mathbb{C}\setminus\mathbb{R}$. Note that since $\phi^{[n]}(z)=z+nb$, one easily sees that $C_{\phi}^nk_w=k_{w-n\bar{b}}$ for $w\in\mathbb{C}$. So any function $f\in B^2_\sigma$ orthogonal to the span of the orbit of $k_w$ must vanish at $w_n:=w-n\bar{b}$ for all $n\geq 0$. But $(w_n)_{n\geq 0}$ does not satisfy the Blaschke-Type condition required for zero sets of functions in $B^2_\sigma$, that is
\[
\sum_{n\geq 0}\frac{|\mathrm{Im}(w_n)|}{1+|w_n|^2}\simeq \sum_{n\geq 0}\frac{n}{1+n^2}=\infty.
\]
Hence $f\equiv 0$ and span$(C_{\phi}^nk_w)_{n\geq 0}$ is dense in $B^2_\sigma$. So now let $b\in\mathbb{R}$. Since $(\mathfrak{F}k_w)(t)=e^{-i\bar{w}t}$ for $w\in\mathbb{C}$, it is enough to show that every exponential $e^{iwt}$ is a cyclic vector for $\widehat{C}_\phi=M_{e^{ibt}}$. This is equivalent to the completeness of the sequence $(e^{i(bn+w)t})_{n\geq 0}$ in $L^2[-\sigma,\sigma]$. But this sequence is the image of $(e^{ibnt})_{n\geq 0}$ under  $M_{e^{iwt}}$ which has dense range in $L^2[-\sigma,\sigma]$. Therefore it is sufficient to prove the completeness of $(e^{ibnt})_{n\geq 0}$ in $L^2[-\sigma,\sigma]$. For this we use a classical result of Carleman (see \cite[page 97]{young}): If $(\lambda_n)_{n\geq 0}$ is a sequence of positive real numbers and 
\[
\liminf_{n\to\infty}\frac{n}{\lambda_n}>\frac{\sigma}{\pi},
\]
then $(e^{i\lambda_nt})_{n\geq0}$ is complete in $\mathcal{C}[-\sigma,\sigma]$ (space of continuous functions). This condition is clearly satisfied for $\lambda_n:=bn$ when $0<b<\pi/\sigma$. The case $-\pi/\sigma<b<0$ follows since now $(e^{-ibnt})_{n\geq0}$ is complete and then conjugating. If $|b|=\pi/\sigma$, then the exponentials $(e^{ibnt})_{n\geq0}$ are $2\sigma$-periodic as we saw in the proof of Theorem \ref{cyclic}. It follows that $(e^{ibnt})_{n\in\mathbb{Z}}$ is an orthogonal basis for $L^2[-\sigma,\sigma]$. So $(e^{ibnt})_{n\geq0}$ in not complete in $L^2[-\sigma,\sigma]$ and neither is its image under the injective operator $M_{e^{iwt}}$. This completes the proof.
\end{proof}

The last result hints at the possibility of finding deeper connections between  composition operators on model subspaces $K_\Theta$ of $H^2(\mathbb{C}_+)$ and non-harmonic Fourier analysis.
%===================================================================================================
\begin{thebibliography}{99}  
%====================================================================================================

%\bibitem{Boas}  R. P. Boas Jr., \emph{Entire functions.} Academic Press, Inc., New York, 1954. 

%\bibitem{Appel} M. Appel, P. S. Bourdon and J. Thrall, \emph{Norms of composition operators on the
%Hardy space}, Experiment. Math. 5(1996), 111–117.

%\bibitem{Bourdon} P. S. Bourdon and D. Q. Retsek, \emph{Reproducing kernels and norms of
%of composition operators.} Acta. Sci. Math. (Szeged), 127 (2001), pp. 387–394.
		              
\bibitem{Bayart-Matheron} F. Bayart and E. Matheron, \emph{Dynamics of linear operators}. Cambridge Tracts in Math. (179), 2009.

\bibitem{Bayart} F. Bayart and S. Tapia-García, \emph{Cyclicity of composition operators on the Fock space}. J. Operator Theory 92 (2024), no. 2, 549–577. 47B33 (47A16).

\bibitem{Branges} L. de Branges, \emph{Hilbert spaces of entire functions}, Prentice-Hall, Inc., Englewood Cliffs, NJ, 1968.





%\bibitem{Barbara} C. C. Cowen and B. D. MacCluer, \emph{Some problems on composition operators.Studies on composition operators}, 17–25.
%Contemp. Math., 213, American Mathematical Society, Providence, RI, 1998.

\bibitem{chacon2007composition} G.A. Chacón, G.R. Chacón and J. Giménez. \emph{ Composition operators on spaces of entire functions}, Proc. Amer. Math. Soc. 135 (2007), no. 7, 2205–2218.	
		
\bibitem{cowen2019composition} C. C. Cowen and B. D. MacCluer. \emph{Composition operators on spaces of analytic functions.} Studies in Advanced Mathematics. CRC Press, Boca Raton, FL, 1995. xii+388 pp.	
		
		
%\bibitem{Elliot} S. J. Elliott and M. T. Jury, \emph{Composition operators on Hardy spaces of a half-plane.} Bull. Lond. Math. Soc. 44 (2012), no. 3, 489-495.
		

    \bibitem{Gallardo} Eva A. Gallardo-Gutiérrez and Alfonso Montes-Rodríguez. \emph{The Role of the Spectrum in the Cyclic Behavior of Composition Operators.} Mem. Amer. Math. Soc. 167 (2004), no. 791, viii+81 pp.                                
		
\bibitem{Garcia 1} S. R. Garcia and M. Putinar. \emph{Complex symmetric operators and applications.} Trans. Amer. Math. Soc. 358 (2006), no. 3, 1285–1315. 
		
\bibitem{Garcia 2} S. R. Garcia and M. Putinar. \emph{Complex symmetric operators and applications. II.} Trans. Amer. Math. Soc. 359 (2007), no. 8, 3913–3931.
		
%\bibitem {GP3}S. R. Garcia and M. Putinar, \emph{Interpolation and complex symmetry.} Tohoku Math. J. (2) 60 (2008), no. 3, 423-440.
		
		
\bibitem {Garcia 3}S. R. Garcia and W. R. Wogen, \emph{Complex symmetric partial isometries.} J. Funct. Anal. 257 (2009), no. 4, 1251-1260.
	
	
\bibitem {Garcia 4}S. R. Garcia and W. R. Wogen, \emph{Some new classes of
complex symmetric operators.} Trans. Amer. Math. Soc. 362 (2010), no. 11, 6065--6077.

\bibitem{Linear Chaos} K. Grosse-Erdmann, A. P. Manguillot, {\em  Linear Chaos}, Springer-Verlag London, 2011.

\bibitem{Hilden} H. M. Hilden and L. J. Wallen, \emph{Some cyclic and non-cyclic vectors of certain operators}, Indiana Univ. Math. J. 23 (1974), no. 7, 557–565.

\bibitem{Ikeda} M. Ikeda, I. Ishikawa and S. Yoshihiro, \emph{Boundedness of composition operators on reproducing kernel Hilbert spaces with analytic positive definite functions}, J. Math. Anal. Appl. 511 (2022), no. 1, Paper No. 126048, 30 pp.




        


%\bibitem{matache} V. Matache, \emph{Composition operators on Hardy spaces of a half-plane.} Proc.
%Amer. Math. Soc., 127 (1999), pp. 1483–1491.
		

\bibitem{Nikolski} N. K. Nikolski, \emph{Operators, Functions, and Systems: An Easy Reading: Volume 2: Model Operators and Systems}, Math. Surveys Monogr., vol. 93, Amer. Math. Soc., Providence, RI, 2002.
        
\bibitem{Noor} S. W. Noor and O. R. Severiano. \emph{Complex symmetry and cyclicity of composition operators on $H^2(\mathbb{C}_+)$.} Proc. Amer. Math. Soc. 148 (2020), no. 6, 2469–2476.


        
\bibitem{Ross} W. T. Ross and W. R. Wogen, \emph{Common Cyclic Vectors for Normal Operators}, Indiana Univ. Math. J. 53 (2004), no. 6, 1537–1550.
   

\bibitem{young} R. M. Young, \emph{An introduction to nonharmonic Fourier series}, Academic Press, Inc., San Diego, CA, 2001. xiv+234 pp.
		
		
\end{thebibliography}
	
%----------------------------------------------------------------------------------------------------
	
	
	
\end{document}
